% --- BLOQUE DE PREÁMBULO UNIVERSAL ---
\documentclass[11pt, letterpaper, twoside]{article}

\usepackage[top=2cm, bottom=2.5cm, left=2cm, right=2cm]{geometry}
\usepackage{fontspec}

\usepackage[spanish, bidi=basic, provide=*]{babel}
\babelprovide[import, onchar=ids fonts]{spanish}
\babelprovide[import, onchar=ids fonts]{english}

% Tipografía serif para imitar el Misal Romano
\babelfont{rm}{Noto Serif}

\usepackage{enumitem}
\setlist[itemize]{label=-}

% --- PAQUETES EDITORIALES LITÚRGICOS ---
\usepackage{xcolor}
\usepackage{multicol}
\usepackage{lettrine}
\usepackage{titlesec}
\usepackage{ragged2e}

\definecolor{rubricred}{HTML}{C8102E}

\newcommand{\rubrica}[1]{\textcolor{rubricred}{\small\textit{#1}}}
\newcommand{\vers}{\textcolor{rubricred}{\textbf{V.\ }}}
\newcommand{\resp}{\textcolor{rubricred}{\textbf{R.\ }}}
\newcommand{\cruz}{\textcolor{rubricred}{✠}}

\titleformat{\section}{\Large\bfseries\centering\color{rubricred}}{}{0pt}{\uppercase}
\titleformat{\subsection}{\large\bfseries\centering\color{rubricred}}{}{0pt}{}
\titleformat{\subsubsection}{\normalsize\bfseries\color{rubricred}}{}{0pt}{}

\renewcommand{\LettrineFontHook}{\color{rubricred}}
\setcounter{DefaultLines}{2}
\setlength{\DefaultFindent}{0.2em}
\setlength{\DefaultNindent}{0em}

\setlength{\parindent}{0pt}
\setlength{\parskip}{0.5em}

\begin{document}

% --- Configuración de Seguridad para Compilación en Bruto ---
\catcode`\_=12
\def\VAR#1{[Dato Dinámico]}
\def\BLOCK#1{}

\begin{center}
    {\Huge\textbf{Subsidio Litúrgico Diario}\par}
    \vspace{0.3cm}
    {\large\textsc{\VAR{dia_liturgico | escape_tex}}\par}
    {\small \VAR{fecha_texto | escape_tex}\par}
\end{center}

\vspace{0.8cm}

\begin{multicols}{2}
\RaggedRight

% ==========================================
% LAUDES
% ==========================================
\BLOCK{ if laudes }
\section*{LAUDES}

\BLOCK{ if laudes.salmo1 }
\subsection*{SALMO 1}
\noindent\rubrica{Antífona:} \VAR{laudes.salmo1.antifona | escape_tex} \par\vspace{0.1cm}
\VAR{laudes.salmo1.texto | escape_tex | lettrine}
\par\vspace{0.3cm}
\BLOCK{ endif }

\BLOCK{ if laudes.cantico }
\subsection*{CÁNTICO DEL ANTIGUO TESTAMENTO}
\noindent\rubrica{Antífona:} \VAR{laudes.cantico.antifona | escape_tex} \par\vspace{0.1cm}
\VAR{laudes.cantico.texto | escape_tex | lettrine}
\par\vspace{0.3cm}
\BLOCK{ endif }

\BLOCK{ if laudes.salmo2 }
\subsection*{SALMO 2}
\noindent\rubrica{Antífona:} \VAR{laudes.salmo2.antifona | escape_tex} \par\vspace{0.1cm}
\VAR{laudes.salmo2.texto | escape_tex | lettrine}
\par\vspace{0.3cm}
\BLOCK{ endif }

\BLOCK{ if laudes.evangelio }
\subsection*{LECTURA BREVE Y CÁNTICO}
\VAR{laudes.evangelio.texto | escape_tex | lettrine}
\par\vspace{0.3cm}
\BLOCK{ endif }
\BLOCK{ endif }

% ==========================================
% SANTA MISA 
% ==========================================
\BLOCK{ if misa }
\section*{SANTA MISA}

\BLOCK{ if misa.canto_entrada }
\subsection*{CANTO DE ENTRADA}
\noindent\rubrica{\VAR{misa.canto_entrada.titulo | escape_tex}} \par\vspace{0.1cm}
\VAR{misa.canto_entrada.letra | format_canto}
\par\vspace{0.3cm}
\BLOCK{ endif }

\BLOCK{ if misa.antifona_entrada }
\noindent\rubrica{Antífona de entrada} \par
\VAR{misa.antifona_entrada | escape_tex | lettrine}
\BLOCK{ endif }

\BLOCK{ if misa.gloria }
\subsection*{GLORIA}
\noindent\rubrica{Se dice el Gloria.} \par\vspace{0.3cm}
\BLOCK{ endif }

\BLOCK{ if misa.oracion_colecta }
\subsection*{ORACIÓN COLECTA}
\VAR{misa.oracion_colecta | escape_tex | lettrine}
\BLOCK{ endif }

\BLOCK{ if misa.primera_lectura }
\subsection*{LITURGIA DE LA PALABRA}
\noindent\rubrica{Lectura del \VAR{misa.primera_lectura.cita | escape_tex}.} \par\vspace{0.2cm}
\VAR{misa.primera_lectura.texto | escape_tex | lettrine} \par\vspace{0.2cm}
\vers Palabra de Dios. \par
\resp Te alabamos, Señor.
\BLOCK{ endif }

\BLOCK{ if misa.salmo }
\subsection*{SALMO RESPONSORIAL}
\noindent\rubrica{\VAR{misa.salmo.cita | escape_tex}} \par\vspace{0.1cm}
\resp \VAR{misa.salmo.respuesta | escape_tex} \par\vspace{0.2cm}
\VAR{misa.salmo.texto | escape_tex | lettrine}
\BLOCK{ endif }

\BLOCK{ if misa.segunda_lectura }
\subsection*{SEGUNDA LECTURA}
\noindent\rubrica{Lectura de la \VAR{misa.segunda_lectura.cita | escape_tex}.} \par\vspace{0.2cm}
\VAR{misa.segunda_lectura.texto | escape_tex | lettrine} \par\vspace{0.2cm}
\vers Palabra de Dios. \par
\resp Te alabamos, Señor.
\BLOCK{ endif }

\BLOCK{ if misa.secuencia }
\subsection*{SECUENCIA}
\VAR{misa.secuencia | escape_tex | lettrine} \par\vspace{0.3cm}
\BLOCK{ endif }

\BLOCK{ if misa.aclamacion_evangelio }
\subsection*{ACLAMACIÓN ANTES DEL EVANGELIO}
\resp Aleluya, aleluya. \par
\VAR{misa.aclamacion_evangelio | escape_tex} \par
\resp Aleluya. \par\vspace{0.3cm}
\BLOCK{ endif }

\BLOCK{ if misa.evangelio }
\subsection*{EVANGELIO}
\noindent\rubrica{Lectura del \VAR{misa.evangelio.cita | escape_tex}.} \par\vspace{0.2cm}
\VAR{misa.evangelio.texto | escape_tex | lettrine} \par\vspace{0.2cm}
\vers Palabra del Señor. \par
\resp Gloria a ti, Señor Jesús.
\BLOCK{ endif }

\BLOCK{ if misa.credo }
\subsection*{CREDO}
\noindent\rubrica{Se dice el Credo.} \par\vspace{0.3cm}
\BLOCK{ endif }

\BLOCK{ if misa.oracion_fieles }
\subsection*{ORACIÓN DE LOS FIELES}
\VAR{misa.oracion_fieles | escape_tex} \par\vspace{0.3cm}
\BLOCK{ endif }

\BLOCK{ endif }

% ==========================================
% VÍSPERAS
% ==========================================
\BLOCK{ if visperas }
\section*{VÍSPERAS}

\BLOCK{ if visperas.salmo1 }
\subsection*{SALMO 1}
\noindent\rubrica{Antífona:} \VAR{visperas.salmo1.antifona | escape_tex} \par\vspace{0.1cm}
\VAR{visperas.salmo1.texto | escape_tex | lettrine}
\par\vspace{0.3cm}
\BLOCK{ endif }

\BLOCK{ if visperas.cantico }
\subsection*{CÁNTICO DEL NUEVO TESTAMENTO}
\noindent\rubrica{Antífona:} \VAR{visperas.cantico.antifona | escape_tex} \par\vspace{0.1cm}
\VAR{visperas.cantico.texto | escape_tex | lettrine}
\par\vspace{0.3cm}
\BLOCK{ endif }
\BLOCK{ endif }

\end{multicols}
\end{document}
